\documentclass[11pt]{article}
\usepackage{amsmath,amsthm,amssymb}
\usepackage{hyperref}
\usepackage[margin=1in]{geometry}

\newtheorem{theorem}{Theorem}
\newtheorem{proposition}{Proposition}
\newtheorem{definition}{Definition}
\newtheorem{remark}{Remark}

\title{Claim Promotion Under Incomplete Evidence:\\
Identifiability Limits Beyond Closed-Set Risk Certificates}
\author{Roman Adriashkin\\Independent Researcher\\CogniPrint Research Initiative}
\date{Theory draft v0.1 --- November 2026}

\begin{document}
\maketitle

\begin{abstract}
Finite-sample selective-risk methods can certify a prediction rule relative to
a specified target distribution, but such a certificate need not identify a
different provenance, open-set, or semantic claim. We formalize this gap using
evidence channels and observational equivalence. If two admissible worlds
induce the same evidence-channel law while a claim-relevant target functional
differs, then the target is not identified by that evidence; no evidence-only
test can distinguish the worlds, even with arbitrarily many i.i.d. samples.
We give a closed-channel camouflage construction, a finite Blackwell-style
garbling example, and a claim-poset formulation of admissibility. The results
are a theory draft: their relationship to existing identification, comparison
of experiments, and OOD-impossibility literature remains under external
review, and no novelty claim is made here.
\end{abstract}

\section{Setup}

Let $\Theta$ be a latent source-state space and let a world $D$ be a
probability law over $\mathcal X\times\Theta$. An evidence channel is a
measurable map $e:\mathcal X\to\mathcal Z$. For a channel family $E$, write
$\mu_D^E$ for its induced observation law.

A claim event is a set $C\subseteq\Theta$. Define the world-level target
\[
\eta_C(D)=\Pr_D(\Theta\in C).
\]
For a fixed acceptance rule $g$, define when well-posed
\[
\eta_{C,g}(D)=\Pr_D(\Theta\in C\mid g(X)=1).
\]

\begin{definition}[Evidence identifiability]
A target functional $\eta$ is identified by evidence $E$ over a world class
$\mathcal W$ if
\[
\mu_D^E=\mu_{D'}^E
\quad\Longrightarrow\quad
\eta(D)=\eta(D')
\]
for all $D,D'\in\mathcal W$.
\end{definition}

\section{Observational-equivalence impossibility}

\begin{theorem}[Claim non-identifiability from observational equivalence]
Suppose $D_0,D_1\in\mathcal W$ satisfy
\[
\mu_{D_0}^E=\mu_{D_1}^E
\]
but $\eta(D_0)\neq\eta(D_1)$. Then $\eta$ is not identified by $E$.
Moreover, for every $n$, the $n$-sample evidence laws coincide, and under
equal prior weights no binary test based only on $E$ can distinguish
$D_0$ from $D_1$ with error below $1/2$.
\end{theorem}

\begin{proof}
Equality of the one-observation laws implies equality of the product laws for
every i.i.d. sample size. Therefore every evidence-measurable statistic has the
same distribution under the two worlds. The target cannot be a function of
the evidence law. The total variation distance between the two evidence
experiments is zero, so the optimal equal-prior binary testing error is
$(1-\mathrm{TV})/2=1/2$.
\end{proof}

\section{Closed-channel non-transfer construction}

Let $e_{\rm closed}(X)$ contain only a closed-set predictor output. Let
$D_0$ be a known-only world and let $g$ depend only on this channel. Construct
$D_\pi$ by mixing a fraction $\pi$ of unknown observations whose
$e_{\rm closed}$ law is identical to the known-channel law.

Then
\[
\mu_{D_\pi}^{e_{\rm closed}}
=
\mu_{D_0}^{e_{\rm closed}}.
\]
The unknown fraction among accepted examples is $\pi$. If every accepted
unknown is an error for the known-family prediction and the known-only
selective risk is $r_0$, then
\[
R_{\rm sel}^{D_\pi}(g)
=
\pi+(1-\pi)r_0.
\]

A finite-sample certificate calibrated for $D_0$ remains a statement about
$D_0$ under its assumptions. It is not a certificate for $D_\pi$ merely
because the closed-channel marginal is unchanged.

\section{Claim promotion principle}

If a proposed claim target $\eta_C$ is not identified by the available
evidence $E$ over the declared world class, then an $E$-measurable risk
certificate cannot by itself identify $C$.

A valid promotion therefore requires at least one of: (i) a defensible
restriction of the world class, (ii) an evidence refinement that breaks the
relevant observational equivalence, or (iii) narrowing the claim to a target
that is identified.

\section{Claim poset}

Claims are ordered by logical implication:
\[
C\preceq C'
\quad\Longleftrightarrow\quad
C'\subseteq C,
\]
so $C'$ is stronger. For example, exact known-family membership is stronger
than known-set membership. Known-set and unknown-set membership are
incompatible, not consecutive levels.

If error tolerances are monotone with logical strength, any certified stronger
claim induces a certified weaker claim. Thus admissible claims form a down-set
of the claim poset. Calling the set a lattice ideal requires additional
closure conditions.

\section{Blackwell connection}

The full observation experiment Blackwell-dominates any deterministic
compression of it because the latter is a garbling. This connection must be
used relative to a specified experiment family.

For a finite example, let $X=\{a,b,c,d\}$. Under known state $K$, put mass
$1/2$ on $a,b$; under unknown state $U$, put mass $1/2$ on $c,d$. Define a
compressed channel by
\[
e(a)=e(c)=z_1,\qquad e(b)=e(d)=z_2.
\]
The compressed laws under $K$ and $U$ are identical, whereas the full supports
are disjoint. Full $X$ identifies $K$ versus $U$ perfectly; the compressed
channel contains no information for the equal-prior binary decision problem.

This example shows that closed-channel compression \emph{can} destroy
claim-relevant information. It does not prove that every real closed-set score
is insufficient for every open-set task.

\section{Relation to prior work}

The underlying ideas of observational equivalence, identification, Blackwell
comparison, and unrestricted OOD impossibility have substantial prior
literature. Likewise, finite-sample selective-risk control is established in
RCPS, Learn-then-Test, Conformal Risk Control, LEC, and related work.

Accordingly, the novelty question is narrower: whether an auditable
claim-admissibility calculus, together with composition and non-promotion
rules across heterogeneous evidence channels, yields a useful new formal
object for AI provenance and safety. This remains unestablished.

\section{Conclusion}

\[
\boxed{
\text{A valid risk certificate for one target does not identify a different
claim without assumptions or evidence that identify that claim.}
}
\]

\end{document}
